\section{Introduction and main results}\label{sec:intro}

The Gaussian decomposition principle used in the factor-recovery argument is the classical theorem of Cram\'er~\cite{Cramer1936}.

The exact chi-square law of the sample variance is one of the classical signatures of Gaussian samples. The unrestricted converse for a fixed sample size $n\ge3$ remains a longstanding characterization problem. Existing results impose additional structure, including symmetry or infinite divisibility; see \cite{Ruben1974,Ruben1978,Bondesson1977,GolikovaKruglov2015,Kruglov2013,EjsmontLehner2017}. We do not claim to resolve that unrestricted problem. The sample variance is a chi-square quadratic form, and spherical averaging of its Laplace transform naturally brings in the symmetrized quantity $M_X(s)M_X(-s)$.

For a centered variance-one random variable $X$, let $M_X(s)=\E e^{sX}$ whenever finite and define
\begin{equation}\label{eq:defect-intro}
R_X(s):=\log M_X(s)+\log M_X(-s)-s^2,
\qquad
\Delta_\rho(X):=\sup_{|s|\le\rho}|R_X(s)|.
\end{equation}
We write $d_{\mathrm K}(X,Z)$ for the Kolmogorov distance between the laws of $X$ and $Z$,
\[
d_{\mathrm K}(X,Z):=\sup_x\left|\Pp(X\le x)-\Pp(Z\le x)\right|.
\]
If $X'$ is an independent copy of $X$, then $M_{X-X'}(s)=M_X(s)M_X(-s)$. The difficulty is that reflection removes the Fourier phase. Our stability theorem shows that the factor can nevertheless be recovered, at the optimal logarithmic scale, from local real-axis information alone.

A useful self-improvement is that the reflected defect already supplies the exponential integrability needed later. Indeed, if $\E X=0$, then Jensen's inequality gives $M_X(s)\ge1$. Hence, whenever $M_X$ is finite on $[-\rho,\rho]$ and $\Delta_\rho(X)\le\Delta$,
\begin{equation}\label{eq:auto-envelope-intro}
M_X(\rho)+M_X(-\rho)\le2e^{\rho^2+\Delta},
\qquad
\E e^{\rho|X|}\le2e^{\rho^2+\Delta}.
\end{equation}
Thus no independent exponential-moment assumption is needed in the reflection theorem.

The proof uses two scales. A fixed complex neighborhood yields high-order moment information for $X-X'$; after truncation, this produces a Gaussian product approximation on a disk of radius of order $\sqrt m$. Zero-freeness on this disk allows one to take a logarithm and recover the factor. The exact theorem, by contrast, follows from equality in Jensen's inequality after spherical averaging.

\subsection*{Exact rigidity}
The positive-semidefinite quadratic-form theorem below motivates the reflection condition.
\begin{theorem}[Exact PSD rigidity]\label{thm:exact-psd}
Let $X=(X_1,\ldots,X_d)$ have independent coordinates with $\E X_i=0$ and $\operatorname{Var}(X_i)=\sigma_i^2$. Let $A=A^{\mathsf T}\succeq0$ have rank $r\ge1$, and suppose
\[
X^{\mathsf T}AX\sim\chi_r^2.
\]
For every active coordinate $i$ with $A_{ii}>0$, assume that
\[
M_i(s)M_i(-s)\ge e^{\sigma_i^2s^2}
\]
for all $s$ in some neighborhood of the origin. Then every active $X_i$ is Gaussian, $X_i\sim\Normal(0,\sigma_i^2)$.
\end{theorem}

No separate MGF-existence hypothesis is required in \cref{thm:exact-psd}: the chi-square law of the quadratic form itself implies a square-exponential moment for every active coordinate, and therefore finiteness of its MGF on the whole real line; see \cref{sec:exact}.

Characterizations based on quadratic forms have a long history. Converse results for sample variance and related forms use additional structural assumptions; see, for example, Christoph--Prohorov--Ulyanov, Kruglov, and Golikova--Kruglov \cite{ChristophProhorovUlyanov2001,Kruglov2013,GolikovaKruglov2015}. The condition here is different: infinite divisibility implies the reflected-MGF dominance through the L\'evy--Khintchine formula, but the converse fails and no symmetry is required. For an infinitely divisible centered law with finite variance, the representation gives, for every $s$ such that both $M_X(s)$ and $M_X(-s)$ are finite,
\[
\log M_X(s)+\log M_X(-s)-\sigma^2s^2
 =2\int\left(\cosh(sx)-1-\frac{s^2x^2}{2}\right)\nu(dx)\ge0.
\]
The converse fails. If $X$ takes the values $3$ and $-1/3$ with probabilities $1/10$ and $9/10$, then it is centered and has variance one. It is not infinitely divisible: if a bounded non-degenerate law of support diameter $d$ were an $n$-fold convolution of an i.i.d. root, that root would have support diameter at most $d/n$, and Popoviciu's inequality would give $\operatorname{Var}(X)/n\le d^2/(4n^2)$, which is impossible for large $n$. Its fourth cumulant is $46/9$, so $R_X(s)=23s^4/54+O(s^6)$ and the local dominance holds. This example is asymmetric. In contrast, a Rademacher variable $\varepsilon$ satisfies $R_\varepsilon(s)=2\log\cosh s-s^2=-s^4/6+O(s^6)$, which is negative for all sufficiently small nonzero $s$.

The L\'evy--Khintchine representation used in this comparison is standard; see Sato~\cite{Sato1999}.

\subsection*{Reflection stability}
The quantitative result recovers a factor from local reflected-MGF information without assuming the dominance sign.
\begin{theorem}[Sharp reflection stability]\label{thm:sharp-upper}
Fix $\rho>0$. There exist constants $C_\rho<\infty$ and $\Delta_0(\rho)\in(0,e^{-e})$ such that the following holds. Let $X$ satisfy
\[
\E X=0,\qquad \E X^2=1,
\qquad M_X(s)<\infty\quad(|s|\le\rho),
\]
and suppose
\[
|R_X(s)|\le\Delta,
\qquad |s|\le\rho,
\]
for some $0<\Delta\le\Delta_0(\rho)$. Then
\[
\dK(X,\Normal(0,1))
\le C_\rho\sqrt{\frac{\log\log(1/\Delta)}{\log(1/\Delta)}}.
\]
\end{theorem}

Classical Cram\'er stability begins with a convolution that is already close to a Gaussian convolution in a global probability metric. Sapogov's work is an early cornerstone of this theory~\cite{Sapogov1951}; for the i.i.d. equal-summand case substantially stronger estimates are available, for example in Bobkov, Chistyakov, and G\"otze~\cite{BobkovChistyakovGotze2013}. In the present problem the input is different: $M_X(s)M_X(-s)$ is controlled only on a fixed real interval, and on the Fourier axis the reflected factorization becomes $|\varphi_X|^2$, which contains no phase information. The main task is therefore to convert local transform control into a global statement about one factor.

One could first try to convert the local reflected-MGF control into a global approximation for $X-X'$ and then invoke a general stability theorem for Cram\'er decomposition. Such an argument discards the structure responsible for the sharp modulus. Our proof instead keeps the reflected factorization throughout: high-order moment information constructs a zero-free disk for a truncated factor itself, after which the logarithm of that factor is controlled directly. The logarithmic coefficient step is a finite-disk estimate in the spirit of quantitative Marcinkiewicz theory; related connections between zero geometry and Gaussian approximation appear in Eremenko--Fryntov, Dinh--Ghosh--Tran--Tran, and Michelen--Sahasrabudhe~\cite{EremenkoFryntov2021,DinhGhoshTranTran2025,MichelenSahasrabudhe2026}. The distinctive step here is that the growing zero-free disk is derived from the local reflected-MGF input rather than imposed as a hypothesis.

The order in \cref{thm:sharp-upper} cannot be improved uniformly. In fact, sharpness already occurs in a symmetric class for which no Fourier phase has to be reconstructed.
\begin{theorem}[Symmetric Laguerre lower bound]\label{thm:sharpness}
There are constants $\rho,c,c_-,c_+>0$ and centered variance-one random variables $W_m$ with symmetric densities $f_m$ such that, for all sufficiently large $m$,
\[
c_-\phi(x)\le f_m(x)\le c_+\phi(x),
\qquad
R_{W_m}(s)\ge0\quad(|s|\le\rho),
\]
and, with
\[
\Delta_m:=\sup_{|s|\le\rho}R_{W_m}(s)\to0,
\]
\[
\log\frac1{\Delta_m}=m\log m+O(m),
\qquad
\dK(W_m,\Normal(0,1))\ge c m^{-1/2}.
\]
Consequently, for some $c'>0$,
\[
\dK(W_m,\Normal(0,1))
\ge c'\sqrt{\frac{\log\log(1/\Delta_m)}{\log(1/\Delta_m)}}
\]
for all sufficiently large $m$.
\end{theorem}

Thus the logarithmic instability is already present when $M_X(-s)=M_X(s)$ and the factor on the real axis can be recovered by a positive square root. Reflection-induced phase loss is an additional difficulty for nonsymmetric factor recovery, but it is not the sole source of the optimal modulus.
